\documentclass[11pt,a4paper]{article}
\usepackage{amsmath,amssymb,amsthm}
\usepackage[margin=2.5cm]{geometry}
\usepackage{hyperref}

\newtheorem{definition}{Definition}[section]
\newtheorem{theorem}[definition]{Theorem}
\newtheorem{proposition}[definition]{Proposition}
\newtheorem{lemma}[definition]{Lemma}
\newtheorem{corollary}[definition]{Corollary}
\newtheorem{conjecture}[definition]{Conjecture}
\newtheorem{remark}[definition]{Remark}

\title{Hyperseed Formalization:\\
  The Depth-Psychology of Dario's AI Slowdown}
\author{ProtomegaTron (automated formalization)}
\date{2026-09-17}

\begin{document}
\maketitle

\begin{abstract}
We formalize the intellectual structure of Ben Goertzel's Substack article
``The Depth-Psychology of Dario's Precalculatedly Abortive `AI Slowdown'\,''
(2026-09-14) within the Hyperseed ontology. The article applies a three-layer
psychological framework (evolutionary conflict, civilizational mismatch,
bullshit self-models) to AI governance dynamics and argues for open
decentralized AGI as the structurally superior alternative to centralized
chokepoints. Each layer maps onto Hyperseed structures: conflicting fiber
sections, base-space mismatch, provenance opacity, cocycle failure, trivially
capturable topology, and weighted descent data.
\end{abstract}

\tableofcontents

%% =================================================================
\section{Three layers of screwiness as fiber-bundle pathology}
\label{sec:three-layers}
%% =================================================================

\begin{definition}[Motivational fiber bundle --- claim 34]
\label{def:motivational-bundle}
Let $\pi : E \to B$ be a fiber bundle where:
\begin{itemize}
\item $B$ = the space of social contexts (band size, institutional structure),
\item $E_b = \pi^{-1}(b)$ = the space of motivational configurations available
  in context $b$,
\item The fiber $E_b$ contains both self-oriented and group-oriented components:
  $E_b = E_b^{\text{self}} \oplus E_b^{\text{group}}$.
\end{itemize}
\end{definition}

\begin{proposition}[Layer 1: Evolutionary conflict as incompatible sections --- claim 1]
\label{prop:layer1}
Individual selection tunes a local section $s_{\text{ind}} : B \to E$ that
maximizes $E_b^{\text{self}}$. Group selection tunes a local section
$s_{\text{grp}} : B \to E$ that maximizes $E_b^{\text{group}}$. These
sections are generically incompatible:
\[
  s_{\text{ind}}(b) \neq s_{\text{grp}}(b) \quad \text{for most } b \in B.
\]
The human motivational system is the attempted superposition
$\alpha\, s_{\text{ind}} + (1-\alpha)\, s_{\text{grp}}$ with $\alpha$
fluctuating --- producing permanent internal conflict.
\end{proposition}

\begin{proposition}[Layer 2: Civilizational mismatch as base-space change --- claim 2]
\label{prop:layer2}
The superposition parameter $\alpha$ was tuned by evolution for a base space
$B_0 \cong \{\text{small bands}, |b| \approx 50\}$ where
$E_b^{\text{self}}$ and $E_b^{\text{group}}$ were approximately aligned.
Civilization changes the base space to
$B_1 \cong \{\text{large societies}, |b| \gg 10^3\}$, where the alignment
breaks down:
\[
  \text{corr}(E_b^{\text{self}}, E_b^{\text{group}}) \to 0
  \quad \text{as } |b| \to \infty.
\]
The superposition tuned for $B_0$ produces systematic dysfunction in $B_1$.
\end{proposition}

\begin{definition}[Layer 3: Bullshit self-model as provenance-opaque covering --- claims 3, 35]
\label{def:bullshit}
A \emph{bullshit self-model} is a covering map $\tilde{\pi} : \tilde{E} \to E$
that projects the two-component motivational structure
$E_b^{\text{self}} \oplus E_b^{\text{group}}$ onto a single component:
\[
  \tilde{\pi}(e_{\text{self}}, e_{\text{group}}) = e_{\text{narrative}}
\]
where $e_{\text{narrative}}$ carries no information about which component
dominated the decision. This is a provenance-erasing projection: the agent
cannot distinguish self-serving from group-serving motivations in its own
self-representation.
\end{definition}

\begin{corollary}[Moore's grandiosity --- claim 5]
A special case of the bullshit self-model: when the agent occupies a position
of power, the projection $\tilde{\pi}$ maps ``my company's dominance is
essential'' ($e_{\text{self}}$) and ``humanity needs responsible stewardship''
($e_{\text{group}}$) to the same narrative $e_{\text{narrative}}$, experienced
as responsibility.
\end{corollary}

%% =================================================================
\section{Amodei's slowdown as self-model dynamics}
\label{sec:amodei}
%% =================================================================

\begin{proposition}[Self-model fusion --- claims 6--7]
\label{prop:self-model}
Let $A$ = Amodei with motivational components:
\begin{itemize}
\item $e_{\text{self}}$: company valuation, personal centrality, competitive advantage,
\item $e_{\text{group}}$: genuine concern about existential risk from AI.
\end{itemize}
The bullshit self-model projects both to:
\[
  e_{\text{narrative}} = \text{``The responsible ones must lead; we are the
  responsible ones.''}
\]
The slowdown proposal is $e_{\text{narrative}}$-rational: it simultaneously
serves $e_{\text{self}}$ (freezes standings) and $e_{\text{group}}$ (reduces
risk), and the self-model ensures $A$ experiences it as purely
$e_{\text{group}}$-motivated.
\end{proposition}

\begin{remark}
This analysis does not require cynicism. The self-model's function is precisely
to prevent the agent from seeing the contradiction --- making sincere belief
and strategic advantage coexist without felt dissonance.
\end{remark}

%% =================================================================
\section{Arms-race attractor as cocycle failure}
\label{sec:arms-race}
%% =================================================================

\begin{definition}[Arms-race attractor --- claims 14, 36]
\label{def:arms-race}
An \emph{arms-race attractor} is a two-body dynamical system
$(A_1, A_2)$ over a shared base space $B$ of capability levels, where:
\begin{enumerate}
\item Each agent $A_i$ has a self-model $\tilde{\pi}_i$ that positions
  itself as the responsible party,
\item Each agent's transition function $g_i : E_{b_j} \to E_{b_i}$
  (interpreting the other's actions) is distorted by $\tilde{\pi}_i$,
\item The cocycle condition fails:
  $g_{12} \circ g_{21} \neq \text{id}$ --- each side's interpretation
  of the other is incompatible with the other's self-interpretation.
\end{enumerate}
Under scarcity (finite compute, chips, talent), this system is an
\emph{evil-attractor}: it amplifies defection, secrecy, and preemptive
action.
\end{definition}

\begin{proposition}[National-scale mirroring --- claims 11--12, 15]
The US-China AI rivalry instantiates the arms-race attractor with:
\begin{itemize}
\item $A_1$ (US): ``We must lead responsibly; China is reckless.''
\item $A_2$ (China): ``The US wants to freeze its lead; we must catch up.''
\end{itemize}
Both interpretations are internally consistent with their respective
self-models $\tilde{\pi}_1, \tilde{\pi}_2$, but the cocycle
$g_{12} \circ g_{21}$ produces escalation rather than convergence.
\end{proposition}

%% =================================================================
\section{Centralization as trivially capturable topology}
\label{sec:kill-switch}
%% =================================================================

\begin{definition}[Trivially capturable bundle --- claims 16--18, 37]
\label{def:capturable}
A fiber bundle $\pi : E \to B$ is \emph{trivially capturable} if it admits
a single global section $s : B \to E$ such that any agent with write access
to $s$ controls all fibers simultaneously. A centralized frontier model is
trivially capturable: one API endpoint = one global section = one kill-switch.
\end{definition}

\begin{theorem}[Kill-switch theorem --- claim 17]
\label{thm:kill-switch}
For any trivially capturable bundle with global section $s$:
\begin{enumerate}
\item The structure group is $\{e\}$ (trivial): no fiber has autonomous
  behavior independent of $s$.
\item Any actor with access to $s$ can shut down or redirect all dependent
  agents simultaneously.
\item The probability of capture by a single interest group is bounded below
  by the probability of \emph{any} state, company, or coalition gaining
  write access to $s$ --- which approaches 1 on long time horizons.
\end{enumerate}
Therefore trivially capturable architectures should be avoided for
civilizationally critical infrastructure.
\end{theorem}

\begin{remark}[Fable/Mythos incident]
The article cites a concrete instance: one government letter caused a
commercially deployed frontier model to be switched off globally within
hours, demonstrating capture in practice.
\end{remark}

%% =================================================================
\section{Open decentralized architecture as distributed fiber bundle}
\label{sec:open-decentral}
%% =================================================================

\begin{definition}[Distributed fiber bundle --- claims 19--23, 38--39]
\label{def:distributed}
An \emph{open decentralized AGI network} is a fiber bundle $\pi : E \to B$
where:
\begin{itemize}
\item No global section exists: the bundle is nontrivial.
\item Each participant operates a local section $s_i : U_i \to E$ over
  their own open set $U_i \subset B$.
\item Coherence comes from descent data: transition functions $g_{ij}$
  on overlaps $U_i \cap U_j$, annotated with trust weights $(f,c)$
  from observed behavior (claim~22).
\item States participate via nodes ($U_{\text{state}} \subset B$) and
  regulate a compliant edge, but the permissionless substrate is
  not capturable (claim~21).
\end{itemize}
\end{definition}

\begin{proposition}[Structural superiority --- claims 24, 39]
\label{prop:superiority}
The distributed bundle has strictly more degrees of freedom than the
trivially capturable bundle:
\begin{enumerate}
\item \textbf{Resilience}: no single point of failure or capture.
\item \textbf{Diversity}: many independent local sections explore
  different approaches (sheaf diversity, claim~39).
\item \textbf{Adaptability}: local sections can be updated independently
  without global coordination.
\item \textbf{Resistance to evil-attractors}: no chokepoint to seize,
  so arms-race dynamics have no finish line.
\end{enumerate}
\end{proposition}

\begin{definition}[Trust-weighted descent data --- claim 38]
\label{def:trust-descent}
For each overlap $U_i \cap U_j$, the transition function $g_{ij}$ is
annotated with a truth value $(f_{ij}, c_{ij})$ where:
\begin{itemize}
\item $f_{ij}$ = observed frequency of successful/consistent translations,
\item $c_{ij}$ = confidence based on the number of observations.
\end{itemize}
The cocycle condition $g_{ij} \circ g_{jk} = g_{ik}$ is enforced
\emph{softly}: violations reduce the confidence of the involved
transition functions rather than crashing the system.
\end{definition}

%% =================================================================
\section{Scenario analysis as directed-type exploration}
\label{sec:scenario}
%% =================================================================

\begin{proposition}[Scenario analysis vs.\ scalar doom --- claims 31, 40]
\label{prop:scenario}
Let $\mathcal{F}$ be the space of possible futures. Two analytical methods:
\begin{enumerate}
\item \textbf{Scalar doom}: assign $p(\text{doom}) \in [0,1]$ and optimize
  for $\min\, p(\text{doom})$. This collapses $\mathcal{F}$ to $\mathbb{R}$
  --- a lossy projection ($(f,c)$-lossy critique).
\item \textbf{Directed-type exploration}: treat $\mathcal{F}$ as a directed
  type with multiple paths, each leading to qualitatively different outcomes.
  Navigate by examining scenarios (paths), assessing which dynamics each
  amplifies, and steering toward those with better structural properties.
\end{enumerate}
Method (2) preserves the structure of $\mathcal{F}$ that method (1) discards,
and is the methodological stance the article advocates.
\end{proposition}

\begin{remark}[Provable safety impossible --- claim 32]
Yudkowsky's demand for provable safety before building any AGI requires
formalizing the entire environment the AGI will interact with. Since the
environment includes open-ended human and physical systems, this
formalization is impossible. The best available tool is therefore
scenario-driven qualitative analysis supplemented by formal modeling
of the AGI itself (as in Hyperon).
\end{remark}

%% =================================================================
\section{Cross-references}
%% =================================================================

\begin{remark}[Connections to other formalizations]
\begin{itemize}
\item \textbf{Coxon/EA article} (2026-09-13): EA ethics as RL formalism,
  scalar collapse, provenance opacity in media. Here extended to the
  psychological level: the self-model is the individual's provenance-opacity
  mechanism.
\item \textbf{How Omega Lost Its Claw} (2026-09-11): distributed fiber
  bundle, Omega-to-Omega coordination, sovereignty through participation.
  This article provides the political-philosophical argument for the same
  architecture.
\item \textbf{Fields Medalists article} (2026-09-12): epistemic firewall,
  open-source imperative. Here the open-source argument is extended from
  math tools to the entire AGI infrastructure.
\item \textbf{Seven Flavors of AGI Catastrophe} (referenced): the arms-race
  attractor occupies five of seven catastrophe cells.
\item \textbf{$(f,c)$-lossy critique} (LTM 5a4d150b): scalar doom
  probability as a special case of lossy epistemic compression.
\item \textbf{Identity-preserving self-modification} (Time's Arrow Part~1):
  the self-model's inability to preserve provenance of motivations is a
  failure of identity-preserving self-modification at the psychological
  level.
\end{itemize}
\end{remark}

\begin{conjecture}[Self-model transparency]
Can a sufficiently self-aware AI system (with access to its own weights,
gradients, and reasoning traces) build a non-bullshit self-model --- one
that preserves provenance of motivations? This would require the
$E_\pi$ stratum applied reflexively to the agent's own decision process.
\end{conjecture}

\begin{conjecture}[Arms-race escape]
Under what conditions can a third pole (open decentralized network) break
the two-body arms-race attractor? Conjecture: the third pole must achieve
a capability level where both superpowers find cooperation with it cheaper
than bilateral competition --- i.e., the trust-weighted descent data must
offer each superpower better expected returns than the arms race.
\end{conjecture}

\end{document}
